% \begin{document}
\section{Duality}
\subsection{Intuition}

Consider the following linear program:
\begin{center}
\begin{tabular}{lll}
  Minimize    & $7x_1+3x_2$ & \\
  Subject to: & $x_1+x_2\ge 2,$  & $/\cdot y_1$\\
              & $3x_1+x_2\ge 4,$ & $/\cdot y_2$\\
              & $x_1,x_2\ge 0.$  &
\end{tabular}
\end{center}

We are looking for the optimal solution $\mathrm{OPT}$ to this LP. To find the solution we may ask
two types of questions (to find the upper and lower bound on $\mathrm{OPT}$). Is there a solution
of cost $\le 10$? (Is $\mathrm{OPT}\le 10$?) An answer to this type of questions is quite simple;
we just find a feasible solution to the LP with objective function $\le 10$ e.g.\ $x_1=x_2=1$. Is
there no better solution? (Is $\mathrm{OPT}\ge 10$?) Observe that we can bound the objective
function value using constraints that every feasible solution satisfies. From the first constraint
we get
\[
  7x_1+3x_2\ge x_1+x_2\ge 2.
\]
Thus $\mathrm{OPT}\ge 2$. Similarly from the second constraint we get $\mathrm{OPT}\ge 4$. To make
a better lower bound we will need to be more clever. Let's take a linear combination of the
constraints with coefficients $y_1,y_2$ correspondingly. $y_1,y_2$ should be non-negative because
multiplying a constraint by negative number would flip the inequality. By taking $y_1=1$, $y_2=2$
we obtain
\[
  7x_1+3x_2\ge(x_1+x_2)+2(3x_1+x_2)\ge 2+2\cdot 4=10.
\]
For each constraint of the given LP we associate a dual variable $y_i$ denoting the weight of the
$i$-th constraint. What kind of variables can we take to get a valid lower bound? How should we
pick coefficients to maximize lower bound for $\mathrm{OPT}$? First of all we are interested in
lower-bounding the objective function. Thus linear combination of primal constraints cannot
exceed primal objective function. As mentioned earlier $y_i\ge 0$. In this way we get the
following (dual) linear program for $y_1,y_2$:
\begin{center}
\begin{tabular}{lll}
  Maximize    & $2y_1+4y_2$ & (lower bound as tight as possible)\\
  Subject to: & $y_1+3y_2\le 7,$ & (coefficient of $x_1$)\\
              & $y_1+y_2\le 3,$  & (coefficient of $x_2$)\\
              & $y_1,y_2\ge 0.$  &
\end{tabular}
\end{center}
Let's try now to formalize this approach.

\subsection{General case}

Consider the following linear program with $n$ variables $x_i$ for $i\in[1,n]$ and $m$
constraints:
\begin{center}
\begin{tabular}{ll}
  Minimize    & $\displaystyle\sum_{i=1}^n c_i x_i$\\[8pt]
  Subject to: & $\displaystyle\sum_{i=1}^n A_{ji}x_i\ge b_j\quad\forall j=1,\ldots,m,$\\[8pt]
              & $x\ge 0.$
\end{tabular}
\end{center}
Then, the dual program has $m$ variables $y_j$ for $j\in[1,m]$ and $n$ constraints:
\begin{center}
\begin{tabular}{ll}
  Maximize    & $\displaystyle\sum_{j=1}^m b_j y_j$\\[8pt]
  Subject to: & $\displaystyle\sum_{j=1}^m A_{ji}y_j\le c_i\quad\forall i=1,\ldots,n,$\\[8pt]
              & $y\ge 0.$
\end{tabular}
\end{center}
Each variable $y_j$ in the dual program corresponds to the weight of one of constraints from the
primal LP. We have $n$ constraints in the dual, one for every primal variable $x_i$.

\paragraph{Remark} We showed how to produce dual program for minimization problem. Similar
approach works also for maximization problems. Also every maximization problem can be reduced to
a minimization one. We replace $x_i$ with $-x_i$ in constraints and objective function obtaining
minimization problem.

\paragraph{Remark} One can verify that if we take the dual of the dual problem, we get back to the
primal problem, as we should expect. Finding the dual linear program is an automatic procedure.

\subsection{Duality Theorems}

Let's focus now on the solutions of both LPs. In our example optimal solutions to primal and dual
problems coincided. We now present two theorems that connect primal and dual solutions.
\label{thm:weak}
\thm{Weak Duality}{
If $x$ is primal-feasible (meaning that $x$ is a feasible solution to the primal problem) and $y$
is dual-feasible, then
\[
  \sum_{i=1}^n c_i x_i\ \ge\ \sum_{j=1}^m b_j y_j.
\]
}

\pf{}{
Let's rewrite the right hand side
\[
  \sum_{j=1}^m b_j y_j\ \le\ \sum_{j=1}^m\sum_{i=1}^n A_{ji}x_i y_j
  \ =\ \sum_{i=1}^n\Big(\sum_{j=1}^m A_{ji}y_j\Big)x_i\ \le\ \sum_{i=1}^n c_i x_i.
\]
Here we used the fact that $x,y\ge 0$ for the inequalities.
}

This theorem tells us that every dual-feasible solution is a lower bound to any primal solution.
This is intuitive: every primal feasible solution satisfies primal constraints, dual feasible
solution gives us a way of lower bounding primal solution using primal constraints. Moreover from
Weak Duality we can conclude that optimal solution to primal program is lower bounded by optimal
solution to dual program. In fact optimal solutions to primal and duals linear programs coincide,
leading to the following theorem.
\label{thm:strong}
\thm{Strong Duality}{
If $x$ is an optimal primal solution and $y$ is an optimal dual solution, then
\[
  \sum_{i=1}^n c_i x_i=\sum_{j=1}^m b_j y_j.
\]
Furthermore, if the primal problem is unbounded, then the dual problem is infeasible and
analogously if the dual is unbounded, the primal is infeasible.
}

We omit the proof of strong duality for now but I encourage you to understand one of the many
proofs!

\subsection{Example: Maximum cardinality matching and Vertex Cover on Bipartite Graphs}

Let $G=(A\cup B,E)$ be a bipartite graph and let $M$ be a matching. Let $x_e$ be a variable
corresponding to taking edge $e\in M$. We want to maximize the cardinality of $M$ while assuring
that every vertex has at most one neighboring edge belonging to $M$. Writing those conditions in a
form of LP gives us:
\begin{center}
\begin{tabular}{ll}
  Maximize    & $\displaystyle\sum_{e\in E}x_e$\\[8pt]
  Subject to: & $\displaystyle\sum_{e=(a,b)\in E}x_e\le 1\quad\forall a\in A,$\\[8pt]
              & $\displaystyle\sum_{e=(a,b)\in E}x_e\le 1\quad\forall b\in B,$\\[8pt]
              & $x_e\ge 0.$
\end{tabular}
\end{center}
Thus the dual program looks like this:
\begin{center}
\begin{tabular}{ll}
  Minimize    & $\displaystyle\sum_{v\in A\cup B}y_v$\\[8pt]
  Subject to: & $y_a+y_b\ge 1\quad\text{for }(a,b)\in E,$\\[4pt]
              & $y_v\ge 0.$
\end{tabular}
\end{center}
One can easily notice that this LP is vertex cover relaxation. By weak-duality, we have that
$|M|\le|C|$ for any matching $M$ and vertex cover $C$. Moreover, since both the primal and the
dual are integral for bipartite graphs, strong LP-duality implies K\"onig's theorem:

\thm{K\"onig 1931}{
Let $M^\star$ be a maximum cardinality matching and $C^\star$ be a minimum vertex cover of a
bipartite graph. Then
\[
  |M^\star|=|C^\star|.
\]
}

Another well-known duality that is also a special case of LP-duality is the max-flow $=$ min-cut
theorem.

\subsection{Complementarity Slackness}

Strong duality gives an important relationship between primal and dual optimal solutions.
\label{thm:cs}
\thm{}{
Let $x\in\mathbb{R}^n$ be a feasible solution to the primal and let $y\in\mathbb{R}^m$ be a
feasible solution to the dual. Then
\[
  x,y \text{ are both optimal solutions}
  \iff
  \begin{cases}
    x_i>0\ \Rightarrow\ c_i=\displaystyle\sum_{j=1}^m A_{ji}y_j & \forall i=1,\ldots,n,\\[10pt]
    y_j>0\ \Rightarrow\ b_j=\displaystyle\sum_{i=1}^n A_{ji}x_i & \forall j=1,\ldots,m.
  \end{cases}
\]
}

\pf{}{
We will apply the strong duality theorem to the weak duality theorem proof.

\medskip
\noindent$\Rightarrow$\quad Let $x$ be the optimal primal solution. From the weak duality theorem
proof, we have that
\begin{equation}\label{eq:chain}
  \sum_{j=1}^m b_j y_j\ \le\ \sum_{j=1}^m\sum_{i=1}^n A_{ji}x_i y_j
  \ =\ \sum_{i=1}^n\Big(\sum_{j=1}^m A_{ji}y_j\Big)x_i\ \le\ \sum_{i=1}^n c_i x_i.
\end{equation}
Here we used the fact that $x,y\ge 0$. On the other hand by the strong duality theorem
\[
  \sum_{j=1}^m b_j y_j=\sum_{i=1}^n c_i x_i.
\]
So in~\eqref{eq:chain} there are equalities everywhere. Thus
\[
  \sum_{i=1}^n c_i x_i=\sum_{i=1}^n\Big(\sum_{j=1}^m A_{ji}y_j\Big)x_i
  \ \Rightarrow\ 
  c_i x_i=\Big(\sum_{j=1}^m A_{ji}y_j\Big)x_i\quad\text{for } i=1,\ldots,n.
\]
And finally for every $x_i$, $i=1,\ldots,n$:
\[
  x_i\neq 0,\ \ c_i x_i=\Big(\sum_{j=1}^m A_{ji}y_j\Big)x_i
  \ \Rightarrow\ 
  c_i=\Big(\sum_{j=1}^m A_{ji}y_j\Big).
\]

\medskip
\noindent$\Leftarrow$\quad Similarly to the previous part we know that:
\begin{align*}
  x_i c_i &= \Big(\sum_{j=1}^m A_{ji}y_j\Big)x_i && \forall i=1,\ldots,n,\\
  y_j b_j &= \Big(\sum_{i=1}^n A_{ji}x_i\Big)y_j && \forall j=1,\ldots,m.
\end{align*}
Thus
\[
  \sum_{j=1}^m b_j y_j=\sum_{j=1}^m\sum_{i=1}^n A_{ji}x_i y_j
  =\sum_{i=1}^n\Big(\sum_{j=1}^m A_{ji}y_j\Big)x_i=\sum_{i=1}^n c_i x_i.
\]
The above equality is equivalent to $x,y$ being optimal solutions to the primal and the dual
linear programs, respectively. Indeed for feasible solution $x^\star$ to the primal we have by
weak duality
\[
  \sum_{i=1}^n c_i x^\star_i\ \ge\ \sum_{j=1}^m b_j y_j=\sum_{i=1}^n c_i x_i.
\]
Thus $x$ is an optimal solution to the primal program and similarly $y$ is an optimal solution to
the dual.
}
% \end{document}
